\begin{proposition}[Reduction of projective variance to separated chains]
\label{prop:projective-variance-reduction}
Fix an odd prime~$p$, and let
\[
 \pi_p\colon\{0,\ldots,p-1\}\longrightarrow
 \mathbf P^1(\mathbf F_p)
\]
be a projective orbit satisfying
\[
 \pi_p(p-1-n)=\pi_p(n)
 \qquad(0\le n<p).
\]
For each $q\in\mathbf P^1(\mathbf F_p)$, let $\Gamma_p(q)$ be any
collection of quadruples
\[
 \gamma=(z_0,z_1,z_2,z_3),
 \qquad 0\le z_0<z_1<z_2<z_3<p,
\]
with the following properties:
\begin{enumerate}
 \item the $z_i$ are four consecutive occurrences of~$q$ in the
 ordered fiber $\pi_p^{-1}(q)$;
 \item membership in $\Gamma_p(q)$ is preserved by
 \[
  R(z_0,z_1,z_2,z_3)
  =(p-1-z_3,p-1-z_2,p-1-z_1,p-1-z_0);
 \]
 \item no member has a reflection-centered adjacent pair, that is,
 $z_i+z_{i+1}\ne p-1$ for $i=0,1,2$.
\end{enumerate}
Put
\[
 C_p(q)=|\Gamma_p(q)|,
 \qquad M_p=\sum_q C_p(q),
 \qquad
 V_p=\sum_q\left(C_p(q)-\frac{M_p}{p+1}\right)^2.
\]
Then $R$ acts freely on every $\Gamma_p(q)$.  If
$\overline\Gamma_p(q)=\Gamma_p(q)/\langle R\rangle$ and
\[
 \overline C_p(q)=|\overline\Gamma_p(q)|,
 \quad \overline M_p=\sum_q\overline C_p(q),
 \quad
 \overline V_p=
 \sum_q\left(\overline C_p(q)-\frac{\overline M_p}{p+1}\right)^2,
\]
then
\[
 C_p(q)=2\overline C_p(q),\qquad
 M_p=2\overline M_p,\qquad
 V_p=4\overline V_p.
\]

Call two distinct reflection orbits \emph{overlapping} if some oriented
representative of one has closed interval $[z_0,z_3]$ meeting the closed
interval of some oriented representative of the other.  Let
\[
 \begin{aligned}
 E_p^{\mathrm{sep}}
 =\#\bigl\{(\overline\gamma,\overline\delta):{}&
  \overline\gamma\ne\overline\delta,\quad
  \overline\gamma,\overline\delta\in\overline\Gamma_p(q)
  \text{ for some }q,\\[-2pt]
 &\overline\gamma,\overline\delta\text{ are not overlapping}
 \bigr\}.
 \end{aligned}
\]
where ordered pairs are counted.  Then
\[
 \sum_q\overline C_p(q)^2
 \le 7\overline M_p+E_p^{\mathrm{sep}},
 \qquad
 \overline V_p\le7\overline M_p+E_p^{\mathrm{sep}}.
\]
If, moreover, a member of $\Gamma_p(q)$ is uniquely determined by its
first index, then $M_p\le p$ and hence $\overline M_p\le p/2$.
Consequently, for every $X\ge2$,
\[
 \sum_{X<p\le2X}\overline V_p
 \ll \frac{X^2}{\log X}
      +\sum_{X<p\le2X}E_p^{\mathrm{sep}}.
\]
In particular, the weak projective-dispersion estimate follows once the
separated energy on the right is
$O(X^{2+o(1)}/\log X)$.
\end{proposition}

\begin{proof}
The reflection identity for~$\pi_p$ shows that $R$ preserves the fiber.
It also preserves consecutiveness: an occurrence of~$q$ strictly between
two reflected indices would reflect to an occurrence strictly between the
corresponding original indices.  Thus condition~2 makes $R$ an involution
of $\Gamma_p(q)$.  If $R\gamma=\gamma$, comparison of the two middle
coordinates gives $z_1=p-1-z_2$, contrary to condition~3.  The action is
therefore free, which proves the three scaling identities, including
\[
 \begin{split}
 V_p
 &=\sum_q\left(2\overline C_p(q)
       -\frac{2\overline M_p}{p+1}\right)^2
 =4\overline V_p.
 \end{split}
\]

Fix a state~$q$, and write its ordered occurrence list as
\[
 r_1<r_2<\cdots<r_t.
\]
Every primitive chain in $\Gamma_p(q)$ is one of the filtered sliding
windows
\[
 \gamma_i=(r_i,r_{i+1},r_{i+2},r_{i+3}).
\]
If $i<j$ and the closed intervals of $\gamma_i$ and $\gamma_j$ meet,
then $r_j\le r_{i+3}$, and hence $j-i\le3$.  The same argument with the
two indices reversed shows that a fixed oriented window overlaps at most
six other oriented windows.  Notice that this remains true after imposing
arbitrary filters: deleting windows cannot create a new overlapping pair.

Choose one oriented representative~$\gamma$ of a fixed reflection orbit
$\overline\gamma$.  If another orbit~$\overline\delta$ overlaps it, then,
after applying~$R$ simultaneously to both representatives when necessary,
one of the two representatives of $\overline\delta$ overlaps~$\gamma$.
The preceding sliding-window bound therefore gives at most six overlapping
orbits $\overline\delta$ for each $\overline\gamma$.  Summing first over
$\overline\gamma$ and then over~$q$, the number of ordered distinct
overlapping orbit pairs is at most $6\overline M_p$.

The square energy $\sum_q\overline C_p(q)^2$ counts the diagonal orbit
pairs, the distinct overlapping pairs, and the distinct separated pairs.
The three contributions are respectively $\overline M_p$, at most
$6\overline M_p$, and $E_p^{\mathrm{sep}}$.  This proves the asserted
energy bound.  Since
\[
 \overline V_p
 =\sum_q\overline C_p(q)^2
  -\frac{\overline M_p^2}{p+1},
\]
the variance bound follows as well.

If the first index determines the chain, the map
$\gamma\mapsto z_0$ injects $\bigcup_q\Gamma_p(q)$ into
$\{0,\ldots,p-1\}$, so $M_p\le p$.  Finally,
\[
 \sum_{X<p\le2X}\overline M_p
 \le X\bigl(\pi(2X)-\pi(X)\bigr)
 \ll\frac{X^2}{\log X}.
\]
Summing the pointwise variance bound completes the proof.
\end{proof}

\begin{corollary}[Distinguished-fiber consequence]
\label{cor:projective-variance-distinguished}
Retain the notation and hypotheses of
Proposition~\ref{prop:projective-variance-reduction}.  Fix
$q_0\in\mathbf P^1(\mathbf F_p)$ and put
$\overline A_p=\overline C_p(q_0)$.  If
\[
 \sum_{X<p\le2X}E_p^{\mathrm{sep}}
 \ll\frac{X^{2+o(1)}}{\log X},
\]
then
\[
 \sum_{X<p\le2X}\overline A_p
 \ll\frac{X^{3/2+o(1)}}{\log X}.
\]
The corresponding oriented count $A_p=C_p(q_0)$ is exactly twice
$\overline A_p$.
\end{corollary}

\begin{proof}
For primes in the dyadic interval, Cauchy--Schwarz gives
\[
 \begin{split}
 \sum_{X<p\le2X}\overline A_p
 &\le
 \sum_{X<p\le2X}\frac{\overline M_p}{p+1}
 +\sum_{X<p\le2X}
  \left|\overline C_p(q_0)-\frac{\overline M_p}{p+1}\right| \\
 &\le
 \sum_{X<p\le2X}\frac{\overline M_p}{p+1}
 +\left(
   \bigl(\pi(2X)-\pi(X)\bigr)
   \sum_{X<p\le2X}\overline V_p
  \right)^{1/2}.
 \end{split}
\]
The first term is $O(X/\log X)$ because
$\overline M_p\le p/2$, while the proposition bounds the second by
$X^{3/2+o(1)}/\log X$.  Finally, $A_p=2\overline A_p$ by freeness of
reflection.
\end{proof}
